\section{Limitations and the Open Challenge}
\label{sec:limitations}

\paragraph{Honest residuals.}
Two corruptions remain unsolved from the membrane alone, and the theory predicts why.
\texttt{spatial\_dropout}'s discriminative signature is \emph{spatial} (which regions went
dark), and global average pooling discards spatial layout before the moments are taken; the
released spatial-dispersion features (\texttt{phi\_spatial}: per-channel spatial variance and
participation ratio) restore that signal and feed the MDD's spatial branch, which we have
implemented and validated end-to-end on a controlled contraction (a signal present only in
\texttt{phi\_spatial} is lifted from fused-without-spatial AUROC $\approx 0.50$ to
$\approx 1.00$). Confirming the lift on real \texttt{spatial\_dropout} requires re-extracting
$\varphi$ with \texttt{phi\_spatial} (a GPU run on the user's cluster) and is the single
highest-leverage remaining experiment. \texttt{polarity\_flip} is a property of the trained
network's polarity invariance, not of our representation, and is undetectable from
$V_{\mathrm{mem}}$ by construction -- its signal lives on the input side. We frame both as an
\textbf{open challenge to the community}: the Achilles' heel of event-based perception, where
corruptions are at once \emph{harmful} to detection (Section~\ref{sec:motivation}) and
\emph{hard} to detect cheaply.
\advisor{point: if we cannot push these above 0.6, frame as a discovery / community alert. Both levers (spatial branch; adaptive neurons for flood) are identified.}

\paragraph{Other limitations.}
(i)~The cost claim is \emph{architectural} -- $V_{\mathrm{mem}}$ is measured in simulation
(SpikingJelly on a GPU), not on neuromorphic silicon; an on-chip energy measurement is future
work. (ii)~AUROC figures are oracle-labelled separability point estimates on a single
leakage-safe split, without bootstrap confidence intervals or a full severity sweep with CIs.
(iii)~Results are on a single dataset (Gen1); transfer to other event datasets (e.g.\ DSEC,
Gen4) is untested. (iv)~Severities are synthetic transforms calibrated to be monotonic, not
measured sensor failures -- the contribution is the benchmark, protocol, and feature release,
not inventing corruption physics.

\section{Conclusion}
\label{sec:conclusion}

We introduced the task of out-of-distribution detection for event-camera spiking detectors
from the \emph{free} sub-threshold membrane potential, the Gen1-C corruption benchmark, a
closed-form theory of per-corruption detectability, and the Manifold-Decomposition Detector
that follows from a dilation/contraction geometric diagnosis. We showed empirically that a
pretrained SOTA detector collapses under these corruptions, motivating cheap on-chip
detection; that $\varphi$ is the only OOD signal that runs inside the target hardware's
compute envelope; and that one unsupervised score takes four of six corruptions to
$\geq 0.96$ with sequence aggregation. The two residuals are predicted by the theory and
released as an open challenge, with a concrete, biologically-grounded path -- spatial-
dispersion features and adaptive neurons -- toward closing them.
